%!TeX root = ../main.tex

\section{Trabalhos futuros}

Este trabalho construiu um ambiente de testes unificado e o aplicou a sete problemas de emparelhamento em grafos. Algumas extensões ficaram fora do escopo, por limitação de tempo ou por decisão metodológica, e constituem os próximos passos naturais:

\begin{itemize}
    \item \textbf{Algoritmos descritos e não implementados:} Parte dos métodos apresentados nos capítulos de problemas não foi incorporada ao ambiente de testes: o algoritmo de Dinic, para cardinalidade máxima em grafos bipartidos; a variante de caminhos mínimos sucessivos com Dijkstra e potenciais, para a atribuição linear; o método dual, para a atribuição gargalo; o algoritmo de Micali e Vazirani, para cardinalidade máxima em grafos gerais; e a versão otimizada do algoritmo de Edmonds ponderado \cite{galil_micali_gabow}. Como o ambiente já define a interface padrão e o protocolo de cada problema, incluí-los exige apenas um novo adaptador por algoritmo.
    \item \textbf{Heurísticas para os demais problemas:} A comparação entre métodos exatos e heurísticos foi feita para cardinalidade máxima em grafos bipartidos, peso máximo em grafos gerais e emparelhamento induzido máximo. Ficaram de fora as heurísticas gulosas descritas para a atribuição linear e a atribuição gargalo, a gulosa ordenada para peso máximo em grafos gerais, as heurísticas gulosa, de Karp-Sipser e de crescimento de caminho para cardinalidade máxima em grafos gerais, e as metaheurísticas de busca local para o emparelhamento induzido \cite{furst}.
    \item \textbf{Contagem de operações elementares:} Avaliou-se instrumentar os algoritmos com contadores de operações, como caminhos aumentantes percorridos e arestas examinadas, uma métrica independente de linguagem e de máquina. A ideia foi descartada porque exigiria manter uma cópia modificada de cada repositório ao lado do original intacto. Uma alternativa a investigar é a instrumentação externa, por ferramentas de rastreamento da própria linguagem, sem alterar o código de terceiros.
    \item \textbf{Separação do efeito da linguagem:} As implementações interpretadas e compiladas foram mantidas em tabelas separadas, e o expoente empírico permitiu compará-las apenas quanto ao crescimento. Implementar um mesmo algoritmo nas duas linguagens permitiria medir diretamente quanto da diferença de tempo se deve à linguagem, e não ao algoritmo.
    \item \textbf{Instâncias reais:} As instâncias utilizadas são sintéticas, geradas por famílias parametrizadas. Aplicar o ambiente a instâncias provenientes das aplicações citadas na introdução verificaria se as conclusões se mantêm fora dos cenários controlados.
    \item \textbf{Comunicação dos defeitos encontrados:} Os defeitos identificados em implementações públicas do método Húngaro (Seção~\ref{sec:defeitos}) podem ser reportados aos seus mantenedores, junto com as instâncias que os reproduzem.
    \item \textbf{Problemas fora do escopo:} As variações NP-difíceis citadas na introdução, como o emparelhamento tridimensional e o problema da atribuição quadrática, e os emparelhamentos em multigrafos e hipergrafos, excluídos na delimitação do estudo.
\end{itemize}
